\section{The cubic endpoint}
\label{sec:cubic-endpoint}

The conditional transport theorem contains no special feature of cubic
threefolds.  This section proves the endpoint statement needed for the
application.  We retain
Cai's quantum-connection convention and reconstruct the exact block used
below, so the application does not depend on translating a qualitative
summary of his result.

\subsection{The primitive-sixth block}

Let \(X\subset\mathbf P^4\) be a smooth cubic threefold, let \(P\) be its
hyperplane class, and put \(q=u^2\ne0\).  In the basis
\(1,P,P^2,P^3\), Cai writes the small \(z\)-equation as
\begin{equation}\label{eq:cubic-small-z-equation}
 z^2\partial_zS=(K+zG)S,
\end{equation}
where
\begin{equation}\label{eq:cai-cubic-matrices}
 K=2\begin{pmatrix}
 0&6q&0&36q^2\\
 1&0&15q&0\\
 0&1&0&6q\\
 0&0&1&0
 \end{pmatrix},
 \qquad
 G=\operatorname{diag}\!\left(\frac32,\frac12,-\frac12,-\frac32\right).
\end{equation}
See \cite[Section~3]{Cai}.  Over \(\Q(\sqrt3)(u)\), the eigenvalues of \(K\)
are \(6\sqrt3u,-6\sqrt3u,0,0\), and its zero eigenspace is one rank-two
Jordan block.  Solving the first off-block Sylvester equation gives, on this
block,
\begin{equation}\label{eq:cubic-zero-block-data}
 M_1=\begin{pmatrix}-19/18&0\\0&19/18\end{pmatrix},
 \qquad (R_2)_{21}=-16/81,
\end{equation}
when the Jordan link is normalized to one.  The indicial polynomial is
therefore
\begin{equation}\label{eq:cubic-indicial-polynomial}
 (\rho+19/18)(\rho-1/18)+16/81
   =\rho^2+\rho+\frac5{36}.
\end{equation}
Its roots are \(-1/6\) and \(-5/6\), hence the formal-monodromy eigenvalues
of this block are \(\zeta_6^{-1}\) and \(\zeta_6\).  This direct calculation
recovers the small-connection part of \cite[Proposition~6]{Cai}.  Cai's
flatness argument in the same proposition gauges the big connection back to
the small one and preserves this rank-two block.

\subsection{The point coefficient}

The preceding calculation finds the packet but does not say whether the
point row detects it.  That coefficient follows from one scalar period.  Write
\begin{equation}\label{eq:cubic-quantum-period}
 F(t)=\sum_{d\geq0}\frac{(3d)!}{(d!)^5}t^d
 ={}_2F_3\!\left(
   \begin{matrix}1/3,\,2/3\\1,\,1,\,1\end{matrix};27t
 \right),
\end{equation}
the equality of the two expressions being the identity
\((3d)!=27^d\,d!\,(1/3)_d(2/3)_d\).

\begin{lemma}[Central charge of the point class]
\label{lem:cubic-central-charge}
Let \(X\subset\mathbf P^4\) be a smooth cubic threefold and \(p\in X\).  In
the small quantum limit there are a nonzero constant \(C\) and an integer
\(k\), both independent of \(z\) and fixed by the chosen Gamma and
graph-space normalizations, such that the component of the Gamma point
section along the unit is
\begin{equation}\label{eq:cubic-point-central-charge}
 Z_p(z):=(s_X(\OO_p),\mathbf 1)_X
 =C\,z^{\,k-3/2}F(q/z^2).
\end{equation}
\end{lemma}

\begin{proof}
Read \eqref{eq:gamma-framed-section} from the right.  For a point on a smooth
threefold \(\operatorname{ch}(\OO_p)=[p]\) with no lower-degree terms, and
\([p]\) has cohomological degree \(6\), so
\((2\pi i)^{\deg/2}\operatorname{ch}(\OO_p)=(2\pi i)^3[p]\).  Every
positive-degree term of \(\Ghat_X\) annihilates the top class, so
\(\Ghat_X[p]=[p]\), and \(c_1(X)\cup[p]=0\) gives \(z^{c_1(X)}[p]=[p]\).  The
grading operator satisfies \(\mu[p]=(3-\tfrac32)[p]\), whence
\(z^{-\mu}[p]=z^{-3/2}[p]\).  Collecting these scalars,
\[
 s_X(\OO_p)=(2\pi)^{-3/2}(2\pi i)^3z^{-3/2}\,L_X(0,z)[p].
\]
Pairing against the unit and expanding the fundamental solution in
descendants leaves the one-point generating function with a point insertion,
\(\sum_dq^d\langle[p]/(z-\psi)\rangle_{0,1,d}\), which is the coefficient of
\(\mathbf 1\) in the small \(J\)-function.

That coefficient is \eqref{eq:cubic-quantum-period}.  Givental's mirror
theorem for toric complete intersections gives the small \(I\)-function of a
degree-three hypersurface in \(\mathbf P^4\),
\begin{equation}\label{eq:cubic-i-function}
 I(q,z)=z\sum_{d\geq0}q^d
 \frac{\prod_{m=1}^{3d}(3P+mz)}{\prod_{m=1}^{d}(P+mz)^5},
\end{equation}
and identifies it with the small \(J\)-function up to a mirror map
\cite{GiventalToricCI,CoatesGivental}.  The degree-\(d\) term here carries
\(z^{3d-5d}=z^{-2d}\), so \(I=z\mathbf 1+O(z^{-1})\) and the mirror map is
trivial.  Setting \(P=0\) in \eqref{eq:cubic-i-function} leaves
\[
 z\sum_{d\geq0}q^d\frac{(3d)!\,z^{3d}}{(d!)^5z^{5d}}
 =z\,F(q/z^2),
\]
which is that coefficient up to the power of \(z\) carried by the chosen
normalization.  The accumulated factors---\((2\pi)^{-3/2}(2\pi i)^3\), the
sign convention relating \(L_X(0,z)\) to \(J\), and the overall power of
\(z\) in the \(I\)- and \(J\)-function conventions---are a nonzero constant
times an integral power of \(z\).
\end{proof}

Only the fractional part of the exponent of \(z\) is used below, so the
integer \(k\) in \eqref{eq:cubic-point-central-charge} never enters, and
neither does a convention which replaces \(z\) by \(-z\): that moves the
branch phases but not the exponents.

The two asymptotic regimes of \eqref{eq:cubic-point-central-charge} match the
two parts of \eqref{eq:cai-cubic-matrices}.  The eigenvalues
\(\pm6\sqrt3u\) of \(K\) produce the exponentially large and small terms of
the expansion, while the rank-two zero block---the one carrying the
primitive-sixth formal monodromy---is regular singular and produces the
algebraic terms.  The packet in question is therefore read from the algebraic
part alone.

That part is a Barnes expansion.  For \({}_pF_q\) with \(p\leq q\) and
\(a_k\) distinct modulo the integers, its algebraic terms are
\begin{equation}\label{eq:barnes-algebraic-part}
 \frac{\prod_{j}\Gamma(b_j)}{\prod_{j}\Gamma(a_j)}
 \sum_{k}
 \frac{\Gamma(a_k)\prod_{j\neq k}\Gamma(a_j-a_k)}
      {\prod_{j}\Gamma(b_j-a_k)}\,
 (-x)^{-a_k}\bigl(1+O(x^{-1})\bigr)
\end{equation}
\cite[Equations~16.11.2 and 16.11.8]{DLMF}.  Here \(a=(1/3,2/3)\) and
\(b=(1,1,1)\); the difference \(1/3\) is not an integer, so no logarithms
occur and the two algebraic series are distinct.  The prefactor is
\(1/(\Gamma(1/3)\Gamma(2/3))\), and \eqref{eq:barnes-algebraic-part} gives the
leading coefficients
\begin{equation}\label{eq:cubic-barnes-coefficients}
\begin{aligned}
 C_{1/3}&=\frac{\Gamma(1/3)\Gamma(1/3)}{\Gamma(2/3)^3}
   \cdot\frac{1}{\Gamma(1/3)\Gamma(2/3)}
  =\frac{\Gamma(1/3)}{\Gamma(2/3)^4},\\
 C_{2/3}&=\frac{\Gamma(2/3)\Gamma(-1/3)}{\Gamma(1/3)^3}
   \cdot\frac{1}{\Gamma(1/3)\Gamma(2/3)}
  =\frac{\Gamma(-1/3)}{\Gamma(1/3)^4}.
\end{aligned}
\end{equation}
The Gamma function has no zeros, and \(-1/3\) is not one of its poles, so both
coefficients are nonzero; the branch factors \((-x)^{-1/3}\) and
\((-x)^{-2/3}\) are nonzero on every sector.  Substituting \(x=27q/z^2\) in
\eqref{eq:cubic-point-central-charge} turns them into the powers
\begin{equation}\label{eq:cubic-barnes-z-powers}
 z^{-3/2}x^{-1/3}\sim z^{-5/6},
 \qquad
 z^{-3/2}x^{-2/3}\sim z^{-1/6},
\end{equation}
whose exponents reduce modulo the integers to the roots \(-5/6\) and
\(-1/6\) of \eqref{eq:cubic-indicial-polynomial}, with monodromy factors
\(\zetaSix\) and \(\zetaSix^{-1}\).  Thus the Gamma point section has a
nonzero coefficient along each of the two formal lines of that block.

It remains to pass from the section to the row.  The row is
\(v\mapsto[v,s_X(\OO_p))\), and \eqref{eq:flat-euler-pairing} is
nondegenerate; because of the half-turn \(z\mapsto e^{-\pi i}z\) in its first
slot, it pairs a formal line of monodromy \(\lambda\) with the line of
monodromy \(\lambda^{-1}\).  A nonzero coefficient of \(s_X(\OO_p)\) along
one of the two lines therefore makes the row nonzero on the other.  Both
coefficients are nonzero here, so the conclusion is the same for both packets
and does not depend on resolving which line is paired with which.

\begin{proposition}[Cubic endpoint lemma]
\label{prop:cubic-endpoint}
For every smooth cubic threefold \(X\),
\begin{equation}\label{eq:cubic-point-packet}
 \rrow_X|_{\cP_{\zeta_6}(X)}\ne0,
 \qquad
 \rrow_X|_{\cP_{\zeta_6^{-1}}(X)}\ne0.
\end{equation}
The packet assertion follows from Cai's connection, while the point-row
nonvanishing is the direct Barnes calculation
\eqref{eq:cubic-quantum-period}--\eqref{eq:cubic-barnes-z-powers}.
Both assertions are made at the formal-primary level: the sector-dependent
Barnes phases prove nonvanishing of the intrinsic formal primary projections,
and no comparison between alternative sectorial lifts is used in the endpoint
contrast.  The
formal-primary Boolean is the one recovered from the cyclic row by
Lemma~\ref{lem:cyclic-row-support}.  The
paper uses half-Tate-normalized labels throughout.  Relative to unnormalized
monodromy, the common shift relabels both endpoints equally, so the normalized
primitive-sixth contrast in \eqref{eq:cubic-point-packet} is unchanged.
\end{proposition}

\begin{remark}[Exact regression]
The paper's verification package uses and verifies an explicit basis of
\(\operatorname{im}K^2\oplus\ker K^2\), normalizes the nilpotent link, solves
the first off-block Sylvester equation, and reconstructs
\eqref{eq:cubic-zero-block-data}--\eqref{eq:cubic-indicial-polynomial} at
three nonzero rational specializations of the quantum parameter by exact
rational arithmetic.  It also checks the
hypergeometric recurrence behind \eqref{eq:cubic-quantum-period} and the
projective-space grading
calculation below.  The exact source, checker, and certificate are included
in the archived release at \url{https://doi.org/10.5281/zenodo.21937490}; the
artifact regresses the matrix-to-block and indicial arithmetic at three
rational specializations, the hypergeometric recurrence, the Barnes non-pole
arguments, and the projective-space grading identity.  It does not verify the
Barnes asymptotic theorem or quantum K\"unneth.  The proof remains the
calculation displayed above.
\end{remark}

\subsection{The two endpoints}

For \(\mathbf P^m\), write \(n=m+1\).  At a nonzero quantum parameter
\(q_{\mathbf P}\),
multiplication by \(c_1(\mathbf P^m)\) is a weighted cyclic shift multiplied
by \(n\).  After adjoining the \(n\)-th root of the parameter and rescaling
the basis, it is \(nq_{\mathbf P}^{1/n}\) times the standard cyclic shift.  In its
Fourier eigenbasis the diagonal of the grading operator is
\begin{equation}\label{eq:projective-grading-average}
 \frac1n\sum_{k=0}^{m}(m/2-k)=0.
\end{equation}
Indeed, if \(F\) is the unitary discrete Fourier matrix, then every diagonal
entry of the Fourier conjugate of the grading operator is
\(\sum_j |F_{jk}|^2(m/2-j)=n^{-1}\sum_j(m/2-j)\), since
\(|F_{jk}|^2=1/n\).  Thus the vanishing in
\eqref{eq:projective-grading-average} is entrywise, not merely a trace
identity.
All \(n\) rank-one formal factors therefore have formal-monodromy exponent
zero.  In particular, \(\mathbf P^{m+3}\) has no primitive-sixth packet.

The quantum connection of a product is the tensor connection.  Hence
\(X\times\mathbf P^m\) has \(m+1\) copies of each cubic primitive-sixth line.
The extended Gamma framing is multiplicative under products
\cite[Remark~1.2]{IritaniGamma}, and the rank row is the tensor product of
the two rank rows.  Since the rank row of \(\mathbf P^m\) is nonzero,
\eqref{eq:cubic-point-packet} implies
\begin{equation}\label{eq:stabilized-endpoint-contrast}
 \rrow_{X\times\mathbf P^m}
   |_{\cP_{\zeta_6}(X\times\mathbf P^m)}\ne0,
 \qquad
 \cP_{\zeta_6}(\mathbf P^{m+3})=0.
\end{equation}

\begin{proof}[Proof of Theorem~\ref{thm:intro-cubic-conditional}]
If \(X\times\mathbf P^m\) were rational, it would be birational to
\(\mathbf P^{m+3}\).  Theorem~\ref{thm:birational-point-primary} would identify
the two point-row Booleans at \(\zeta_6\).
Equation~\eqref{eq:stabilized-endpoint-contrast} gives the values
one and zero, a contradiction.
\end{proof}

The separation between Theorems~\ref{thm:birational-point-primary} and
\ref{thm:intro-cubic-conditional} is deliberate.  The conditional transport
argument does not use Cai's calculation.  Any replacement endpoint with a
point-row-detected primary packet gives another application of the same
criterion.
